\documentclass[11pt,a4paper,addpoints]{exam}
\usepackage[utf8]{inputenc}
\usepackage[T1]{fontenc}
\usepackage[polish]{babel}
\usepackage{geometry}

\pointpoints{pkt}{pkt}
\boxedpoints
\hqword{Zadanie:}
\hpword{Punkty:}
\hsword{Wynik:}
\htword{Suma}

\lhead{\makebox[12cm][l]{\large\textbf{Imię i nazwisko:}\enspace\hrulefill}}
\rhead{W01: git}

\footer{}{}{}

\begin{document}

\cellwidth{2.75cm}
\noindent \gradetable[h]

\vspace{0.8cm}

\begin{questions}

\question[2] 
Rozpoczynasz pracę nad nowym zadaniem, które składa się z~trzech etapów: (1)~Utworzenie struktury, (2)~Dodanie logiki, (3)~Napisanie testów. Przyjmij, że wszystkie zmiany w kolejnych etapach dotyczą wyłącznie pliku Program.cs. Podaj kolejno komendy, których użyjesz, aby utworzyć lokalne repozytorium, zacommitować każdy z~trzech etapów osobno, a~następnie przesłać całość do~nowego, wcześniej utworzonego, pustego repozytorium na~GitHubie.
\fillwithdottedlines{6cm}

\question[2] 
W Twoim repozytorium istnieje główna gałąź \texttt{main}. Od niej utworzono poboczną gałąź o~nazwie \texttt{feature-auth}, na której zapisałeś 2~commity. Zakończyłeś pracę. Jesteś obecnie na gałęzi \texttt{feature-auth}. Podaj komendy, przy pomocy których zintegrujesz ukończone zmiany z~gałęzi \texttt{feature-auth} do gałęzi \texttt{main}. Załóż, że wystąpił konflikt scalania.
\fillwithdottedlines{6cm}

\newpage

\question[1] 
Pracując nad nową funkcjonalnością na swojej gałęzi, zapisałeś postępy w ostatnich~5~małych commitach. Przed wysłaniem kodu na serwer chcesz połączyć te 5~commitów w~jeden, pozostając na tej samej gałęzi. Podaj komendy, dzięki którym zmienisz historię w~ten sposób, nie tracąc przy tym żadnego napisanego kodu.
\fillwithdottedlines{6cm}

\end{questions}

\end{document}
